% Artículo VII. Incorporación de resultados preexistentes, 10-09-2026.
% Propietarios íntegros preservados en fuentes_conservadas:
% 17_pantalla_cubica_soldadura.tex (P), 18_clausura_energetica.tex (E).
% P: incidencia, cociente, soldadura celular y refinamiento.
% E: respuesta, unicidad, extensión mínima y expectativa tracial.
% Dependencias anteriores de VII: núcleo APP--TRIT--TPK; realización
% cúbica de la ruta, cociclo de transporte y compresión q_vis=5/9.
% Este fragmento no constituye por sí solo el artículo autónomo.
% La identidad de pantalla se distingue de la contracción material de espín.

\section{Soudure de frontière et réponse énergétique}
\label{vii:sec:pantalla-respuesta}

La réalisation cubique du transport conserve six incidences orientées.
Son quotient géométrique, la forme lorentzienne et la réponse énergétique
s’obtiennent à partir de la même matrice d’incidence. Le transport provient du
chemin enrichi APP--TRIT--TPK : la face, le signe, la retenue et la
mémoire accompagnent chaque arête. Cette section développe
l’opération linéaire qui agit sur cette représentation de frontière.

\subsection{Incidence cubique et quotient de rang quatre}
\label{vii:sec:cociente-pantalla}

Soit \(H_F=\mathbb R^F\), avec
\(F=\{(i,\epsilon):1\leq i\leq3,\ \epsilon\in\{-1,1\}\}\),
son produit scalaire euclidien et sa base \(e_{i,\epsilon}\).
L'opposition des faces est l'involution orthogonale
\(Re_{i,\epsilon}=e_{i,-\epsilon}\). Nous définissons
\[
 u=\sum_{i,\epsilon}e_{i,\epsilon},\qquad
 P_u=\frac{uu^{\mathsf T}}6,\qquad P_-=\frac{I-R}{2},\qquad
 P_0=\frac{I+R}{2}-P_u,\qquad P_\square=P_u+P_-.
\]
Les sommets sont \(V=\{-1,1\}^3\). La matrice d'incidence \(M:\mathbb R^V\longrightarrow H_F\) est déterminée par
\[
 M_{(i,\epsilon),\nu}=\mathbf1_{\{\nu_i=\epsilon\}}.
\]

\begin{proposition}[Décomposition de l'incidence]
\label{vii:prop:incidencia}
On a
\[
 MM^{\mathsf T}=12P_u+4P_-,
 \qquad
 \ker M^{\mathsf T}=P_0H_F.
\]
Par conséquent, le quotient
\(Q=H_F/P_0H_F\), identifié orthogonalement à
\(P_\square H_F\), est de dimension quatre.
\end{proposition}

\begin{proof}
Une face contient quatre sommets. Deux faces opposées n'ont pas de sommets
communs et deux faces d'axes différents ont deux sommets communs.
En écrivant \(\mathbf J=uu^{\mathsf T}\), on obtient
\[
 MM^{\mathsf T}=4I+2(\mathbf J-I-R)
 =2(I-R)+2\mathbf J=4P_-+12P_u.
\]
Les trois projecteurs \(P_u,P_-,P_0\) sont orthogonaux et leur somme vaut \(I\) ; leurs rangs sont respectivement un, trois et deux. L'identité \(\|M^{\mathsf T}x\|^2=\langle x,MM^{\mathsf T}x\rangle\) identifie le noyau annoncé. Le rang complémentaire est quatre.
\end{proof}

Sur \(Q\), la forme
\[
 B=P_--P_u
\]
a pour signature \((3,1)\). L'orientation temporelle est fixée par \(u\).
Avec
\[
 t=\frac{u}{\sqrt6},\qquad
 d_i=\frac{e_{i,+}-e_{i,-}}{\sqrt2},\qquad
 W=\sqrt6P_u+\sqrt2P_-,
\]
on a \(We_{i,\epsilon}=t+\epsilon d_i\). Ces six vecteurs
sont nuls pour \(B\), tandis que
\[
 3t+\nu_1d_1+\nu_2d_2+\nu_3d_3
\]
a pour norme quadratique \(-6\). La même incidence satisfait
\(MM^{\mathsf T}=2W^*W\) sur \(Q\). Les rotations propres du cube
commutent avec les projecteurs, préservent les deux formes et fixent \(t\).
Ces identités déterminent la métrique de la représentation cubique.

\subsection{Soudure cellulaire et transport de mémoire}
\label{vii:sec:soldadura-celular}

Pour chaque arête orientée \(a\), la représentation du parcours fournit la face \(\lambda(\underline a)\), le signe \(\sigma(a)\), les fibres d'origine et d'arrivée et le transport \(U(a):Q_{s(a)}\longrightarrow Q_{t(a)}\). L'inversion vérifie \(U(\bar a)=U(a)^{-1}\). Les demi-arêtes de frontière conservent leur inversion formelle et leur incidence au bord. Avec \(\ell_m=\ell_0\,9^{-m}>0\), la soudure est
\begin{equation}
 \beta_m(a)=\ell_m\sigma_m(a)
              n_{\lambda(\underline a)}^{s(a)},\qquad
 n_{i,\epsilon}=t+\epsilon d_i.
 \label{vii:eq:soldadura-arista}
\end{equation}
La face et le signe sont transportés avec la mémoire de la route. La convention d'inversion exige
\begin{equation}
 \beta_m(\bar a)=-U_m(a)\beta_m(a).
 \label{vii:eq:inversion-soldadura}
\end{equation}
En appliquant deux fois cette relation, on retrouve \(\beta_m(a)\) ; le registre de retenue conserve séparément la composition ordonnée.

\begin{proposition}[Naturalité de la soudure par raffinement]
\label{vii:prop:refinamiento-soldadura}
Soit \(a=a_1\cdots a_9\) un raffinement compatible d’une arête.
Supposons que les neuf incidences filles, transportées vers la fibre du
prédécesseur, conservent la face et le signe, et que
\(\ell_{m+1}=\ell_m/9\). Soit
\(\mathcal U_j:Q_{m,s(a)}\longrightarrow Q_{m+1,s(a_j)}\)
le transport cumulé qui inclut l’identification par raffinement
et le trajet jusqu’à l’origine du successeur \(j\). Alors
\[
 \beta_m(a)=
 \sum_{j=1}^{9}\mathcal U_j^{-1}\beta_{m+1}(a_j).
\]
L'égalité est compatible avec l'inversion et avec la concaténation du registre de mémoire.
\end{proposition}

\begin{proof}
Chaque terme, exprimé dans la fibre de l’origine du prédécesseur, est
\((\ell_m/9)\sigma_m(a)n_{\lambda(\underline a)}\).
La somme de neuf termes reproduit la soudure du prédécesseur.
L’inversion renverse l’ordre des transports et remplace chacun
par son inverse ; la relation \eqref{vii:eq:inversion-soldadura} fournit
le signe global. La mémoire se concatène dans ce même ordre, de sorte
que le raffinement conserve tant la représentation vectorielle que son
registre enrichi.
\end{proof}

\subsection{Réponse positive, unicité et extension minimale}
\label{vii:sec:respuesta-positiva}

La soudure linéaire d'écran est \(\bar\beta_m=\ell_mW|_Q\). Elle est inversible, avec
\[
 \bar\beta_m^{-1}
 =\ell_m^{-1}\left(\frac1{\sqrt6}P_u+\frac1{\sqrt2}P_-\right).
\]
La valeur propre \(5/9\) du secteur antisymétrique du bloc normalisé \(B_c/9\), construit dans la section~\ref{vii:sec:bloque-compresion}, détermine le coefficient de compression \(q_{\mathrm{vis}}=5/9\). Nous réservons cette notation pour le distinguer du quotient euclidien \(q_9(n)\) du noyau arithmétique. Sa réalisation sur l'écran utilise une isométrie \(V_9:Q\longrightarrow\mathcal K\). La contribution visible \(C_9=q_{\mathrm{vis}}V_9\) et son complément \(D_9=\sqrt{1-q_{\mathrm{vis}}^2}\,V_9\) satisfont
\[
 C_9^*C_9+D_9^*D_9=I_Q,\qquad D_9^*D_9=\frac{56}{81}I_Q.
\]
Le coefficient \(56\) est \(9^2-5^2\) ; il appartient au défaut quadratique
de ce lecteur de compression.

\begin{theorem}[Réponse énergétique déterminée par la soudure]
\label{vii:thm:respuesta-unica}
Il existe un unique opérateur \(\Lambda_m:Q\longrightarrow Q\) tel que
\(\bar\beta_m^*\Lambda_m\bar\beta_m=D_9^*D_9\). Il est défini positif et
\begin{equation}
 \Lambda_m=\frac{28}{243\ell_m^2}(P_u+3P_-)
 =\frac{112}{81\ell_m^2}
       (MM^{\mathsf T}|_Q)^{-1}.
 \label{vii:eq:lambda-pantalla}
\end{equation}
Sur \(H_F\), l'identité correspondante a pour second membre
\((56/81)P_\square\).
\end{theorem}

\begin{proof}
La congruence par la soudure inversible détermine nécessairement
\[
 \Lambda_m=\bar\beta_m^{-*}\frac{56}{81}I_Q\bar\beta_m^{-1}.
\]
En substituant l'inverse, le coefficient de \(P_u\) est \(56/(486\ell_m^2)=28/(243\ell_m^2)\) ; celui de \(P_-\) est \(28/(81\ell_m^2)\). Tous deux sont positifs. Par la proposition \ref{vii:prop:incidencia}, \((MM^{\mathsf T}|_Q)^{-1}=P_u/12+P_-/4\), ce qui donne la seconde expression. L'extension à six faces annule \(P_0H_F\) ; l'identité s'étend donc au moyen de \(P_\square\).
\end{proof}

\begin{proposition}[Caractérisation par extension minimale]
\label{vii:prop:minima-extension}
Pour tout \(b\in Q\),
\[
 \langle b,\Lambda_m b\rangle
 =\frac{112}{81\ell_m^2}
   \min_{Mx=b}\|x\|^2.
\]
L'extension minimisante est unique :
\(x_{\min}=M^{\mathsf T}(MM^{\mathsf T}|_Q)^{-1}b\).
\end{proposition}

\begin{proof}
Le vecteur annoncé satisfait \(Mx_{\min}=b\) et appartient à
\(\operatorname{im}M^{\mathsf T}=(\ker M)^\perp\).
Toute autre solution est \(x_{\min}+z\), avec \(z\in\ker M\) ; par
orthogonalité, sa norme au carré est \(\|x_{\min}\|^2+\|z\|^2\).
De plus,
\[
 \|x_{\min}\|^2
 =\langle b,(MM^{\mathsf T}|_Q)^{-1}b\rangle.
\]
La formule de \(\Lambda_m\) conclut la preuve.
\end{proof}

L'opérateur \(MM^{\mathsf T}|_Q\) réalise la réponse
de Dirichlet--à--Neumann de cette formulation d'incidence ; son inverse
est la réponse de Neumann--à--Dirichlet. L'énergie précédente correspond
à cette dernière, avec la normalisation de la soudure et du défaut de
compression. Le complément de Schur de
\(\left(\begin{smallmatrix}I&M^{\mathsf T}\\M&0\end{smallmatrix}\right)\)
est \(-MM^{\mathsf T}\), ce qui offre la même élimination variationnelle.

\subsection{Conservation de la densité par raffinement}
\label{vii:sec:densidad-refinamiento}

Considérons le secteur des transports cubiques, orthogonaux pour le produit scalaire énergétique et compatibles avec la forme \(B\). Nous fixons un espace commun d'étiquettes \(Q^{\mathrm{lab}}\), identifié isométriquement à l'écran du prédécesseur. Les soudures ont les types \(\bar\beta_m:Q^{\mathrm{lab}}\longrightarrow Q_m\) et \(\bar\beta_{m+1,j}:Q^{\mathrm{lab}}\longrightarrow Q_{m+1,j}\) ; les opérateurs \(\Lambda\) agissent dans leurs fibres d'arrivée respectives. Si \(U_j:Q_j\longrightarrow Q\) transporte la fibre du successeur vers celle du prédécesseur, alors
\[
 \bar\beta_{m+1,j}=\frac19U_j^{-1}\bar\beta_m,\qquad
 \Lambda_{m+1,j}=81U_j^{-1}\Lambda_mU_j.
\]
Le facteur \(81\) procède de \(\ell_{m+1}^{-2}=81\ell_m^{-2}\).

\begin{theorem}[Invariance de la densité normalisée]
\label{vii:thm:densidad}
Chaque successeur conserve la réponse quadratique du prédécesseur :
\[
 \bar\beta_{m+1,j}^*\Lambda_{m+1,j}\bar\beta_{m+1,j}
 =\frac{56}{81}I_{Q^{\mathrm{lab}}}.
\]
Pour \(r\) raffinements, la réponse conjointe, exprimée dans les fibres du prédécesseur, est
\[
 E_{m+r}=\frac{56}{81}I_{Q^{\mathrm{lab}}}\otimes I_{9^r}.
\]
L'espérance de trace normalisée
\(\tau_{9^r}(A)=9^{-r}\operatorname{Tr}(A)\) satisfait
\[
 (\operatorname{id}\otimes\tau_{9^r})(E_{m+r})
 =\frac{56}{81}I_{Q^{\mathrm{lab}}}
\]
à toute profondeur finie.
\end{theorem}

\begin{proof}
La congruence d'un successeur contient les facteurs \(1/9,81,1/9\) ; leur produit vaut un. L'orthogonalité annule les transports intermédiaires. La somme directe de neuf réponses égales identifie le premier raffinement à \((56/81)I_Q\otimes I_9\). L'itération produit le tenseur avec \(I_{9^r}\). Enfin \(\tau_{9^r}(I_{9^r})=1\). Les espérances se composent suivant l'identité \(\tau_{9^{r+s}}=\tau_{9^r}\otimes\tau_{9^s}\).
\end{proof}

Les inclusions unitales \(A\mapsto A\otimes I_9\) sont compatibles avec cette collection de traces. Dans la limite inductive des algèbres matricielles, la réponse est l'élément central \((56/81)I_Q\otimes\mathbf1\). La somme extensive non normalisée produit en revanche \(9^r(56/81)I_Q\). Cette distinction spécifie la grandeur conservée : la densité quadratique d'écran.

\subsection{Réalisation variationnelle de la réponse}
\label{vii:sec:realizacion-respuesta}

La réalisation matérielle utilise une isométrie \(J_{\mathrm{scr}}:Q\longrightarrow\mathcal H_\partial\) compatible avec l'action de spin et avec le raffinement. La cotétrade de frontière est \(e=J_{\mathrm{scr}}\bar\beta_m\). Soit \(\Lambda_{\mathrm{CE}}\) la réponse de frontière obtenue par seconde variation de l'action matérielle et géométrique, avec ses conditions aux limites et l'élimination des variables intérieures. Le défaut de compatibilité est
\begin{equation}
 \mathcal R_{\mathrm{grav}}
 =\bar\beta_m^*
   (J_{\mathrm{scr}}^*\Lambda_{\mathrm{CE}}J_{\mathrm{scr}}-\Lambda_m)
   \bar\beta_m
 =e^*\Lambda_{\mathrm{CE}}e-D_9^*D_9.
 \label{vii:eq:defecto-respuesta}
\end{equation}
Par l'inversibilité de \(\bar\beta_m\),
\[
 \mathcal R_{\mathrm{grav}}=0
 \quad\Longleftrightarrow\quad
 J_{\mathrm{scr}}^*\Lambda_{\mathrm{CE}}J_{\mathrm{scr}}=\Lambda_m.
\]
L'isométrie \(J_{\mathrm{scr}}\) transporte les états de frontière ;
le courant variationnel de spin est une forme de degré trois à valeurs
dans les bivecteurs. Chacun intervient avec son domaine propre dans la
réalisation de la réponse.

La construction détermine une forme énergétique positive, son extension intérieure minimisante et la densité conservée par raffinement. L'amplitude d'une contribution matérielle concrète s'obtient en évaluant le courant de cette réalisation dans sa forme quadratique variationnelle. Sont ainsi distingués l'opérateur qui fixe la réponse, le courant sur lequel il agit et l'observable gravitationnel qui exprime l'évaluation.
